\SourceLeaf{359}%
daetische Linie von bekannter Grösse und Richtung verbunden ist, ist schon seit langer Zeit vielfältig behandelt, und um unter verschiedenen Methoden zu seinem Gebrauch passend zu wählen, muss man allerdings mancherlei Umstände berücksichtigen. Es ist z.\,B. erheblich dabei, ob man die Aufgabe nur für Einen oder einige wenige concrete Fälle aufzulösen hat, oder für sehr viele. In der letztern Voraussetzung wird es von Wichtigkeit sein, dass die Methode jedesmal die möglich grösste Bequemlichkeit und Übersichtlichkeit der Definitivrechnung gewähre, wenn auch die Anwendbarkeit der Methode vielleicht erst gewisse allgemeine Vorbereitungsarbeiten erfordern sollte. Eben so wichtig ist der Umstand, ob man die Resultate einer ausgedehnten trigonometrischen Vermessung alle in der Form von geographischer Länge und Breite und zwar ausschliesslich \textit{nur} in \textit{dieser} Form verlange, oder ob daneben die Resultate für die Lage sämmtlicher Punkte auch noch in einer andern Form, z.\,B. der der rechtwinkligen Coordinaten, aufgestellt werden; im letztern Fall wird es weniger nothwendig sein, die geographische Lage mit der alleräussersten Schärfe anzugeben.

Die von \textsc{Dusejour}, \textsc{Legendre}, \textsc{Delambre} u.\,A. gegebenen Formeln berücksichtigen nur die erste Potenz der Abplattung, was allerdings in practischer Hinsicht von nicht grosser Erheblichkeit sein wird, da einmal die Abplattung des Erdsphäroids nur ein kleiner Bruch ist. Es ist daher auch nicht die Meinung, es als einen in practischer Beziehung wichtigen Vorzug geltend zu machen, dass die neue Methode von der Kleinheit der Abplattung ganz unabhängig ist. Die bessern unter jenen Methoden mögen allerdings eine in den meisten Fällen zureichende Schärfe gewähren, obwohl man einen in mathematischer Beziehung genügenden Nachweis dafür vermisst. Dagegen darf man behaupten, dass die neue Methode, wenn die nöthigen Erfordernisse bereit sind, eine bequemere und nach ihrem wesentlichen Inhalt in einem bedeutend kleinern Raum zu concentrirende Rechnung ergibt. \textsc{Bessels} im Jahre 1825 gegebene Auflösung trägt das Gepräge einer grossen mathematischen Vollendung, und ist auch gar nicht abhängig von der Voraussetzung, dass die Entfernung der beiden Punkte von einander im Vergleich zu den Dimensionen des ganzen Erdsphäroids klein sei. In theoretischer Rücksicht ist dies ohne Zweifel ein Vorzug dieser Methode; bei Beurtheilung des practischen Werthes hat man aber folgende Umstände in Betracht zu ziehen. Die Methode macht gar keinen Unterschied zwischen dem Fall grosser und dem Fall kleinerer Entfernungen, sondern erfordert für alle Fälle gleich lange Rechnungen,
